% ---------------------------------------------------------------------------------
% 2. METHODOLOGY
% ---------------------------------------------------------------------------------
\section{Methodology}\label{sec:method}

% 2.1 Guide: pipeline + justification with validation evidence (also what did NOT help).
\subsection{Preprocessing and data augmentation}\label{sec:prep}
Images are read as RGB tensors in $[0,1]$ and normalised with the ImageNet statistics, which the pretrained backbones require; CUSTOM uses the same pipeline. During training, each $320\times320$ image is randomly cropped to $224\times224$, and every batch is then augmented on the GPU with horizontal and vertical flips, $90^\circ$ rotations and random changes of gamma, brightness, contrast and saturation. The augmentation has a \emph{light} and a \emph{strong} preset; the search chose \emph{strong} for CUSTOM and \emph{light} for Swin-T. Validation images are centre-cropped to $224\times224$, the setting used for model selection. For the test predictions, however, we feed the whole $320\times320$ image, which both networks accept because they end in global pooling: although they were trained and validated at $224\times224$, in our experiments the full image gave better results on the test set. No test-time augmentation (TTA) is applied.

An aerial image has no natural ``up'', so flips and $90^\circ$ rotations create new samples without changing the label. Colour is different, since the vegetation is probably the main visual cue for fire risk, so the colour changes were kept mild. We also tried mixup \citep{zhang2018mixup}, CutMix \citep{yun2019cutmix} and TTA, but none of them improved validation macro-F1 by more than the seed noise (\cref{sec:protocol}).

% 2.2 Guide: complete architecture, primitive layers only; design decisions vs reference.
\subsection{System design: CUSTOM network}\label{sec:custom}
The CUSTOM network (\cref{tab:arch}) is built only from primitive layers: four VGG-style blocks \citep{simonyan2015vgg} of 48 to 384 channels, each with two Conv $3\times3$--BatchNorm \citep{ioffe2015batchnorm}--ReLU layers and a $2\times2$ max-pooling, followed by global average pooling (GAP) \citep{lin2014nin} and a small classifier with dropout, for 2.74\,M parameters.

We reached this design step by step from the reference notebook, keeping a change only when it improved validation macro-F1. The first one was to replace the Flatten head, which held almost all the parameters of the network, by GAP. We then used two $3\times3$ convolutions with BatchNorm in each block instead of a single $5\times5$ one, which covers the same area with fewer weights and trained more stably. Finally, since the network was still underfitting, we made it wider (\cref{sec:protocol}).

% 2.3 Guide: architecture, weights, pretraining dataset, new 7-output layer, frozen/updated
% layers, preprocessing matching the weights, per-group learning rates.
\subsection{System design: FINE-TUNING network}\label{sec:ft}
We kept the provided AlexNet \citep{krizhevsky2012alexnet} as a baseline and compared it with three more recent backbones with torchvision ImageNet-1k weights \citep{russakovsky2015imagenet}: ConvNeXt-T \citep{liu2022convnext}, Swin-T \citep{liu2021swin} and EfficientNetV2-S \citep{tan2021efficientnetv2}. EfficientNetV2-S was unstable with mixed precision (its loss became non-finite), so we replaced it by ResNet-50 \citep{he2016resnet}. In every model the original classifier is replaced by a new linear layer with 7 outputs, and the adaptive pooling by a simple mean so that training stays deterministic.

We did not freeze any layer: satellite images are quite different from ImageNet photographs, so the whole network has to adapt. To avoid destroying the pretrained features, the backbone uses a small learning rate and the new head the same rate or ten times more \citep{howard2018ulmfit}, and the inputs follow the pretraining conditions ($224\times224$, ImageNet normalisation). Swin-T obtained the best validation macro-F1 (0.546) and overfitted less than ResNet-50, so it is our FT model.

% 2.4 Guide: loss, optimiser, LR, schedule, batch, epochs, seed, environment, subsampling;
% checkpoint selection, early stopping, reproducibility.
\subsection{Training protocol and model selection}\label{sec:protocol}
\begin{table}[H]
\centering\scriptsize
\setlength{\tabcolsep}{4pt}
\caption{Training configuration of the selected models.}\label{tab:config}
\begin{tabular}{@{}l l l@{}}
\toprule
 & CUSTOM (n07\_60ep) & FT (Swin-T, swin\_p2\_1) \\
\midrule
Loss & \multicolumn{2}{l}{Cross-entropy, class weights $\propto 1/\sqrt{n_c}$, no label smoothing, no mixup} \\
Optimiser \& schedule & \multicolumn{2}{l}{AdamW \citep{loshchilov2019adamw}; linear warm-up + cosine decay \citep{loshchilov2017sgdr}; fp16 mixed precision} \\
Learning rate & $3.17\cdot10^{-4}$ & backbone $2.37\cdot10^{-5}$, head $\times1$ \\
Weight decay & 0.0464 & 0.0395 \\
Batch / warm-up & 128 / 1 epoch & 32 / 2 epochs \\
Epochs (selected epoch) & 60 (36) & 30 (26) \\
Augmentation & strong & light, drop path 0.2 \\
\midrule
Common & \multicolumn{2}{l}{Seed 42; full train split (no subsampling); Kaggle Tesla T4; PyTorch 2.11 \citep{paszke2019pytorch}} \\
\bottomrule
\end{tabular}
\end{table}

\Cref{tab:config} gives the final configuration of both models. To reach it we searched in phases, going from broad to specific.

\textbf{Why random search.} We had about ten hyperparameters to tune (architecture, learning rate, weight decay, batch size, augmentation, class weights, \dots). A grid search would need too many runs even with two or three values per hyperparameter, and most of them would be spent varying hyperparameters that barely matter. For the same budget, a random search tries many more different values of each hyperparameter \citep{bergstra2012random}. We also stopped early the trials that were clearly below the median of the previous ones \citep{golovin2017vizier}, to spend the time on more promising configurations.

\textbf{CUSTOM.} In phases 1 and 2 (about 40 short runs of 15 epochs, split between two computers) we first chose the architecture and the optimiser, and then secondary settings such as the class weights and the regularisation. In phase 3 the three best configurations were trained for 40 epochs, and all of them got stuck at about 0.45--0.46 with almost equal training and validation scores. This told us that the network was underfitting, so instead of tuning further we changed approach: in phases 4 and 5 we tested one idea at a time. Widening the network and training for 60 epochs were the only clear improvements, whereas inverse-frequency class weights made the results worse and more capacity (more channels or blocks) did not help. Since repeating a run with another seed already changes the score by about 0.01, we only accepted changes clearly above that. A final, broader search of 121 trials found a model with a higher validation score, but it did not hold on the public test, so we kept the simpler network.

\textbf{FINE-TUNING.} For each backbone we tried 10 configurations for 10 epochs and retrained the best two for 30 epochs.

In every run we keep the epoch with the best validation macro-F1, without early stopping. Training is deterministic (fixed seeds and deterministic algorithms), so the notebook reproduces the selected models on the same GPU (Tesla T4).
